\section{引言}

\subsection{研究背景与动机}

RTL 设计的安全属性验证要求确认：从合法初始状态出发，设计的执行始终满足给定规格。有界模型检查能够发现有限深度内的可达反例，但未发现反例并不能证明任意长度的执行均安全。归纳证明提供了另一条途径。以直接对目标属性进行一步归纳为例，证明既要确认初始状态满足属性，又要确认任意满足属性的当前状态经过一次转移后仍满足属性。后一条件涵盖实际不可达的状态，因此，即使属性在所有可达状态上成立，它也未必能单独构成归纳不变式。辅助引理通过刻画设计状态之间的约束，能够与目标属性一起构成可归纳的证明条件。构造此类引理通常需要理解 RTL 中的状态更新、目标属性与中间状态关系。

大语言模型（LLM）已被用于提出候选引理，再由形式化工具检查其有效性。例如，Peled 等人的 \emph{Large Lemma Miners}~\cite{peled2026large} 将候选生成与形式化筛选结合起来。候选池的形成并不意味着目标属性已经得到证明：候选可能无法解析、不满足初始条件，或虽正确却与目标证明无关；另一些候选则可能只有联合使用才能发挥作用。因此，候选生成之后仍存在独立的选择问题，即应优先检查哪些引理组合。

这一选择问题同时受验证成本约束。不同组合需要分别接受初始性和归纳步骤检查，失败后的再次尝试还会增加形式化查询；若借助 LLM 分析候选或修复组合，模型推理也会消耗墙钟时间。仅减少求解器调用而不计算新增的模型开销，不能说明整体证明过程更高效。本文因此聚焦于固定候选池中的目标导向选择：在保持最终证明由形式化工具确认的前提下，探索如何用候选间的语义联系和证明反馈安排检查顺序，并控制总成本。

\subsection{研究问题与挑战}

本文固定 RTL 设计、目标安全属性、环境假设及候选引理池；选择阶段只从已有候选中取子集，不生成新引理，也不修改待验证的规格。对于被选组合，初始性要求每个合法初始状态都同时满足目标属性和组合中的每条引理；归纳保持则要求任意满足这些条件的状态经过一步合法转移后，后继状态仍满足它们。只有两项均被形式化后端证明，该组合才能用于证明目标属性。这一定义将候选生成能力与选择算法的效果分开，使不同方法可以在相同候选池和形式化后端上比较。

\emph{Large Lemma Miners}~\cite{peled2026large} 的筛选算法先对候选分类排序，再依次检查排序后的前缀。它只访问有限的前缀组合：如果有效组合不是该顺序的前缀，而夹在其中的候选又妨碍归纳，搜索就可能漏过有效组合。经典 Houdini 则从候选全集出发，迭代删除违反归纳条件的候选；在相关查询均完成且语义一致时，它适合求取固定池的归纳子集，但其搜索过程并非专为尽早证明某一指定属性而安排。两者都提供了重要的比较基准，本文不预设其中任何方法在所有任务上更慢或更弱。

目标导向的组合选择面临三项相互关联的挑战。首先，单条引理的检查结果不能完全说明它在组合中的价值：某条引理的归纳步骤失败，可能意味着它需要其他引理支持，而不是应被永久舍弃。其次，形式化失败必须按初始性、目标属性步骤和辅助引理步骤区分；反例状态若无法按照相同的 SVA、时钟和复位语义可靠求值，就不能用于确定性剪枝。最后，形式化检查、LLM 调用与回退搜索共享同一时间预算，需要在可能的证明收益与实际成本之间作出可审计的决策。

据此，本文围绕以下研究问题展开：

\noindent\textbf{RQ1：}在统一预算下，目标导向的 LLM 引导选择能否在维持有竞争力的证明完成率时，减少形式化查询和包含模型开销的证明时间？\par
\noindent\textbf{RQ2：}当候选池扩大并混入无关、冗余或非归纳候选时，这种选择方式的效率和完成率如何变化？\par
\noindent\textbf{RQ3：}组合级初始选择与结构化证明反馈，能否分别带来超出候选排序和仅接收验证成败的重试的增量收益？

\subsection{主要贡献与论文结构}

为研究上述问题，本文提出 LLM-GLS。该方法让 LLM 基于 RTL、目标属性和固定候选池推荐少量候选 ID 组合及排序。形式化后端逐组检查初始性和归纳性，并把可可靠提取的失败义务反馈给选择器。选择器据此修复组合、过滤被已知反例确定否决的组合；成本控制器再根据剩余预算决定继续检查、请求新组合或进入基于完整合格候选池的 Houdini 回退。LLM 只决定检查的优先级，任何成功输出都必须由形式化证明支持。

本文的主要设计与贡献包括：
\begin{enumerate}
  \item 
  \item 
  \item 
\end{enumerate}

后文先介绍相关工作与形式化问题定义，再通过动机示例展示非前缀组合的必要性；随后给出 LLM-GLS 的算法和系统实现，介绍实验设置与结果分析；最后讨论方法的适用边界与研究局限，并在结论中展望后续工作。
% 正式实验完成后，依据实测结果补充主要定量发现，并复核贡献措辞。
